% =========================================================
% Proof: Monotonicity of Diameter
% Source: volume-iii/analysis/intro-metric-spaces/notes/notes-metric-space.tex
% Mode: standard
% =========================================================

\subsection*{Monotonicity of Diameter}
\label{prf:monotonicity-of-diameter}

\begin{remark*}[Return]
$\leftarrow$ \hyperref[thm:monotonicity-of-diameter]{Back to Theorem
(Monotonicity of Diameter) in Notes}
\end{remark*}

\begin{proof}
\Claim Let $(X, d)$ be a metric space and let $A \subseteq B \subseteq X$.
Then $\operatorname{diam}(A) \le \operatorname{diam}(B)$.

\Given A metric space $(X, d)$ and subsets $A \subseteq B \subseteq X$.
\Goal To show $\operatorname{diam}(A) \le \operatorname{diam}(B)$.
\Strategy We proceed directly by unfolding the definition of diameter and
using the fact that $A \subseteq B$ to bound the supremum over $A$ by
the supremum over $B$.

If $A = \varnothing$ then $\operatorname{diam}(A) = 0 \le \operatorname{diam}(B)$
and the result holds trivially. Assume henceforth that $A \neq \varnothing$.

\ByDef the diameter of a nonempty set $S \subseteq X$ is
\[
  \operatorname{diam}(S) \;=\; \sup\{\, d(x, y) \mid x, y \in S \,\}.
\]

\Let $x, y \in A$ be arbitrary.
Since $A \subseteq B$, we have $x, y \in B$, so $d(x,y)$ belongs to the
set $\{\, d(x,y) \mid x, y \in B \,\}$ over which $\operatorname{diam}(B)$
is taken.
\Therefore $d(x, y) \le \operatorname{diam}(B)$.

Since $x, y \in A$ were arbitrary, $\operatorname{diam}(B)$ is an upper
bound for $\{\, d(x,y) \mid x, y \in A \,\}$.
\ByThm{The supremum is the least upper bound}
$\operatorname{diam}(A) \le \operatorname{diam}(B)$. \AsReq
\end{proof}

\begin{remark*}[Proof shape]
Direct, with a one-line dispatch of the empty case. The subset inclusion
transfers pairs from $A$ into $B$, making $\operatorname{diam}(B)$ an upper
bound on pairwise distances in $A$; the least-upper-bound property of the
supremum then gives the inequality.
\end{remark*}

\begin{remark*}[Dependencies]
\begin{itemize}
  \item Definition of diameter:
        $\operatorname{diam}(S) = \sup\{d(x,y) \mid x, y \in S\}$
        \quad \ref{def:diameter}
  \item Convention: $\operatorname{diam}(\varnothing) = 0$
        \quad \ref{def:diameter}
  \item The supremum is the least upper bound: if $M$ is an upper bound
        for a set $S$ then $\sup S \le M$
        \quad \ref{thm:sup-least-upper-bound}
  \item Definition of subset:
        $A \subseteq B \Rightarrow (x \in A \Rightarrow x \in B)$
        \quad \ref{def:subset}
\end{itemize}
\end{remark*}